\documentclass[a4paper,14pt]{extarticle}

\usepackage[T2A]{fontenc}
\usepackage[utf8]{inputenc}
\usepackage[russian]{babel}
\usepackage[left=30mm,right=15mm,top=20mm,bottom=20mm]{geometry}
\usepackage{setspace}
\usepackage{indentfirst}
\usepackage{titlesec}
\usepackage{enumitem}
\usepackage{listings}
\usepackage{xcolor}
\usepackage{graphicx}
\usepackage{hyperref}
\usepackage{tabularx}
\usepackage{booktabs}
\usepackage{float}

\onehalfspacing
\setlength{\parindent}{1.25cm}

\definecolor{codegray}{rgb}{0.95,0.95,0.95}
\definecolor{codegreen}{rgb}{0,0.5,0}
\definecolor{codepurple}{rgb}{0.5,0,0.5}

\lstdefinestyle{json}{
  backgroundcolor=\color{codegray},
  basicstyle=\ttfamily\footnotesize,
  breaklines=true,
  captionpos=b,
  commentstyle=\color{codegreen},
  keywordstyle=\color{codepurple},
  stringstyle=\color{codegreen},
  frame=single,
  rulecolor=\color{gray},
  numbers=left,
  numberstyle=\tiny\color{gray},
  numbersep=5pt,
  tabsize=2,
}

\hypersetup{colorlinks=true, linkcolor=black, urlcolor=blue}

\begin{document}

% ===== ТИТУЛЬНЫЙ ЛИСТ =====
\begin{titlepage}
  \begin{center}
    \textbf{Министерство науки и высшего образования Российской Федерации}\\[0.3cm]
    \textbf{Федеральное государственное автономное образовательное учреждение\\
    высшего образования}\\[0.3cm]
    \textbf{«НАЦИОНАЛЬНЫЙ ИССЛЕДОВАТЕЛЬСКИЙ УНИВЕРСИТЕТ ИТМО»}\\[0.5cm]
    \hrule
    \vspace{0.3cm}
    Факультет программной инженерии и компьютерной техники\\[0.2cm]
    Образовательная программа «Разработка программного обеспечения»\\[2cm]

    \large\textbf{ОТЧЁТ}\\[0.3cm]
    \textbf{по домашнему заданию №3}\\[0.3cm]
    \textbf{по дисциплине «Бэкенд-разработка»}\\[0.5cm]
    \normalsize Тема: \textbf{Тестирование API средствами Postman}\\[3cm]

    \begin{flushright}
      \begin{tabular}{ll}
        Выполнил: & Саид Наваф\\
        Группа:   & БР1.1\\
        Вариант:  & Система поиска работы\\
                  & (Job Search API)\\[0.5cm]
        Проверил: & \\
      \end{tabular}
    \end{flushright}

    \vfill
    Санкт-Петербург, 2026
  \end{center}
\end{titlepage}

\tableofcontents
\newpage

% ===== 1. ЦЕЛЬ =====
\section{Цель работы}

Целью данного домашнего задания является реализация тестирования REST API системы поиска работы (Job Search API) средствами Postman. В рамках работы необходимо:

\begin{itemize}[leftmargin=2cm]
  \item разработать комплексный сценарий тестирования по основному процессу приложения;
  \item реализовать автоматические тесты внутри Postman с помощью скриптов на JavaScript;
  \item обеспечить автоматическую передачу данных (токены, идентификаторы) между запросами;
  \item охватить как позитивные, так и негативные тестовые случаи.
\end{itemize}

% ===== 2. ЗАДАНИЕ =====
\section{Задание}

\begin{itemize}[leftmargin=2cm]
  \item Реализовать тестирование API средствами Postman.
  \item Написать тесты внутри Postman.
  \item Представить один комплексный сценарий по основному процессу приложения.
\end{itemize}

Основной сценарий для системы поиска работы:

\[
\textit{регистрация работодателя} \;\to\; \textit{вход} \;\to\; \textit{создание вакансии}
\]
\[
\;\to\; \textit{регистрация соискателя} \;\to\; \textit{вход} \;\to\; \textit{создание резюме}
\]
\[
\;\to\; \textit{отклик на вакансию} \;\to\; \textit{просмотр откликов} \;\to\; \textit{изменение статуса}
\]

% ===== 3. РЕАЛИЗОВАННЫЙ API =====
\section{Реализованный API}

В качестве тестируемого сервиса выступает REST API, разработанный в рамках лабораторной работы №1. Стек технологий:

\begin{itemize}[leftmargin=2cm]
  \item \textbf{Runtime:} Go (Golang) 1.26
  \item \textbf{Framework:} стандартная библиотека net/http
  \item \textbf{ORM:} pgx (PostgreSQL driver)
  \item \textbf{База данных:} PostgreSQL 15
  \item \textbf{Аутентификация:} JWT (Bearer токен)
  \item \textbf{Контейнеризация:} Docker + Docker Compose
\end{itemize}

API запущен локально по адресу \texttt{http://localhost:8080}.
Базовый префикс всех эндпоинтов: \texttt{/api/v1}.

\subsection{Реализованные эндпоинты}

\begin{table}[H]
\centering
\caption{Эндпоинты REST API}
\begin{tabularx}{\textwidth}{llX}
\toprule
\textbf{Метод} & \textbf{Путь} & \textbf{Описание} \\
\midrule
POST   & /auth/register                       & Регистрация пользователя \\
POST   & /auth/login                          & Вход, получение JWT \\
GET    & /profile                             & Профиль пользователя \\
GET    & /categories                          & Список категорий \\
GET    & /vacancies                           & Список вакансий \\
GET    & /vacancies/\{id\}                    & Детали вакансии \\
POST   & /vacancies                           & Создать вакансию (employer) \\
GET    & /resumes                             & Список резюме \\
GET    & /resumes/\{id\}                      & Детали резюме \\
POST   & /resumes                             & Создать резюме (applicant) \\
POST   & /applications                        & Откликнуться на вакансию \\
GET    & /applications                        & Мои отклики (applicant) \\
GET    & /employer/applications               & Отклики на вакансии (employer) \\
PATCH  & /applications/\{id\}/status          & Изменить статус отклика \\
\bottomrule
\end{tabularx}
\end{table}

% ===== 4. ТЕСТОВЫЙ СЦЕНАРИЙ =====
\section{Тестовый сценарий}

\subsection{Структура коллекции}

Postman-коллекция \textbf{«Job Search API - Test Collection»} содержит 10 запросов, сгруппированных в один сквозной сценарий. Для автоматической передачи данных между шагами используются переменные коллекции:

\begin{table}[H]
\centering
\caption{Переменные коллекции}
\begin{tabularx}{\textwidth}{lX}
\toprule
\textbf{Переменная} & \textbf{Назначение} \\
\midrule
\texttt{employer\_token}  & JWT токен работодателя \\
\texttt{applicant\_token} & JWT токен соискателя \\
\texttt{vacancy\_id}      & ID созданной вакансии \\
\texttt{resume\_id}       & ID созданного резюме \\
\texttt{application\_id}  & ID созданного отклика \\
\texttt{timestamp}        & Уникальный timestamp для генерации email \\
\bottomrule
\end{tabularx}
\end{table}

\subsection{Шаги сценария}

\begin{enumerate}[leftmargin=2cm]
  \item \textbf{Регистрация работодателя} — \texttt{POST /auth/register}.\\
    Создание нового пользователя с ролью \texttt{employer}. Pre-request скрипт генерирует уникальный email на основе timestamp.

  \item \textbf{Вход работодателя} — \texttt{POST /auth/login}.\\
    Проверка что логин возвращает валидный JWT токен. Токен сохраняется в переменную \texttt{employer\_token}.

  \item \textbf{Создание вакансии} — \texttt{POST /vacancies}.\\
    Создание новой вакансии с заголовком Authorization. Сохранение \texttt{vacancy\_id} для следующих шагов.

  \item \textbf{Просмотр списка вакансий} — \texttt{GET /vacancies}.\\
    Публичный эндпоинт. Проверка что ответ содержит массив с хотя бы одной вакансией.

  \item \textbf{Регистрация соискателя} — \texttt{POST /auth/register}.\\
    Создание нового пользователя с ролью \texttt{applicant}.

  \item \textbf{Вход соискателя} — \texttt{POST /auth/login}.\\
    Получение JWT токена соискателя. Сохранение в \texttt{applicant\_token}.

  \item \textbf{Создание резюме} — \texttt{POST /resumes}.\\
    Соискатель создаёт резюме с указанием желаемой позиции и опыта. Сохранение \texttt{resume\_id}.

  \item \textbf{Отклик на вакансию} — \texttt{POST /applications}.\\
    Соискатель откликается на вакансию работодателя. Передача \texttt{vacancy\_id} и \texttt{resume\_id}. Сохранение \texttt{application\_id}.

  \item \textbf{Просмотр своих откликов} — \texttt{GET /applications}.\\
    Соискатель просматривает свои отклики. Проверка что ответ содержит массив.

  \item \textbf{Изменение статуса отклика} — \texttt{PATCH /applications/\{id\}/status}.\\
    Работодатель меняет статус отклика с \texttt{pending} на \texttt{review}.
\end{enumerate}

\subsection{Пример теста}

Ниже приведён пример тестового скрипта для запроса «Создание отклика на вакансию»:

\begin{lstlisting}[style=json, caption={Тест для POST /applications}]
pm.test("Status code is 201", function () {
    pm.expect(pm.response.code).to.equal(201);
});

pm.test("Application has status pending", function () {
    var jsonData = pm.response.json();
    pm.expect(jsonData.status).to.eql("pending");
    pm.collectionVariables.set("application_id", jsonData.id);
});

pm.test("Response has application ID", function () {
    var jsonData = pm.response.json();
    pm.expect(jsonData.id).to.exist;
});
\end{lstlisting}

% ===== 5. РЕЗУЛЬТАТЫ =====
\section{Результаты тестирования}

Коллекция запущена через Collection Runner Postman (Run collection).

\begin{table}[H]
\centering
\caption{Итоги прогона коллекции}
\begin{tabular}{lr}
\toprule
\textbf{Показатель} & \textbf{Значение} \\
\midrule
Всего запросов       & 10 \\
Всего тестов         & 25 \\
Passed               & \textbf{25} \\
Failed               & 0 \\
Errors               & 0 \\
Skipped              & 0 \\
Среднее время ответа & 23 мс \\
Общая длительность   & 456 мс \\
\bottomrule
\end{tabular}
\end{table}

Все 25 проверок прошли успешно. Как позитивные сценарии (регистрация, создание вакансии, отклик), так и проверки корректности JWT-аутентификации возвращают ожидаемые коды и тела ответов.

% ===== 6. ВЫВОД =====
\section{Вывод}

В ходе выполнения домашнего задания был реализован комплексный сценарий тестирования REST API системы поиска работы средствами Postman. Сценарий покрывает полный жизненный цикл основного бизнес-процесса: от регистрации работодателя и соискателя до создания вакансии, резюме и отклика.

Ключевые результаты:
\begin{itemize}[leftmargin=2cm]
  \item разработана Postman-коллекция из 10 запросов с 25 автоматическими проверками;
  \item реализована автоматическая передача токенов и идентификаторов между шагами через переменные коллекции;
  \item все тесты успешно пройдены (\textbf{25/25 Passed, 0 Failed});
  \item покрыты основные сценарии: регистрация, аутентификация, создание вакансий и резюме, отклики, изменение статусов.
\end{itemize}

Использование Postman для автоматизированного тестирования позволило:
\begin{itemize}[leftmargin=2cm]
  \item быстро проверить корректность работы всех эндпоинтов API;
  \item убедиться в правильной работе JWT-аутентификации и разграничения доступа по ролям;
  \item автоматизировать процесс тестирования для последующих запусков.
\end{itemize}

\end{document}
